\chapter{Introducción}
\label{chap:introducion}

\lettrine{E}{n} este capítulo se exponen las motivaciones que han dado lugar al desarrollo del
proyecto, así como los objetivos concretos que se persiguen. A continuación, se presenta la
solución propuesta para alcanzar dichos objetivos, destacando los principales enfoques
técnicos utilizados. Finalmente, se describe la estructura general de la memoria, detallando
el contenido de cada uno de los capítulos que la componen.

\section{Motivación}
\label{sec:motivacion}

El impulso por la automatización y la precisión en entornos dinámicos es la principal
motivación de este proyecto. Los coches de carreras tipo Scalextric, cuya trayectoria se
restringe a la pista y en los que fundamentalmente hay que controlar la aceleración, son
especialmente aptos para la aplicación de técnicas de visión por computador y de algoritmos
de inteligencia artificial y técnicas de aprendizaje: constituyen un banco de pruebas real,
accesible y controlado en el que experimentar con problemas de percepción y control
en tiempo real.

Siguiendo la línea de trabajos de fin de grado anteriores, se plantea evolucionar el sistema
de seguimiento y control de coches tipo Scalextric presentado el año pasado, que se basaba
en un único nodo con varias cámaras conectadas a través de \gls{usb}, hacia un enfoque donde
el sistema de control lo realice un conjunto de nodos distribuidos conectados por una red
local (wifi y/o Ethernet). El enfoque monolítico presenta una limitación estructural: todas
las cámaras deben conectarse físicamente al mismo equipo, cuyo número de puertos y capacidad
de procesamiento acotan el tamaño del circuito que es posible cubrir. La evolución hacia una
arquitectura distribuida permite desarrollar un sistema más robusto, automático y escalable,
adaptado a circuitos de mayor tamaño y complejidad.

El reto central que aborda este trabajo es estudiar la viabilidad de mantener un control
efectivo y eficaz de los coches en esta arquitectura, de forma que puedan circular a la
máxima velocidad posible en cada tramo del circuito sin salirse en las curvas, teniendo en
cuenta la latencia que introduce la red y el tiempo de procesado de todo el lazo de control:
detección, seguimiento, cálculo de la velocidad máxima en cada tramo y control de la
velocidad.

El desarrollo propuesto puede servir de ayuda en el entrenamiento de las carreras de coches
Scalextric o para determinar el mejor tiempo posible en un circuito. Además, el proyecto
podría iterarse en el futuro para diseñar mecanismos en los que varios coches del mismo
circuito colaboren y aprendan de forma conjunta.

\section{Objetivos del trabajo}
\label{sec:obxectivos}

El objetivo general del proyecto es evolucionar el sistema de control visual de coches tipo
Scalextric hacia una arquitectura distribuida capaz de operar en tiempo real sobre circuitos
de cualquier tamaño. Este objetivo general se concreta en los siguientes objetivos
específicos:

\begin{itemize}
    \item Diseñar una arquitectura modular para el control visual distribuido en tiempo
    real de la velocidad de cada coche en un circuito tipo Scalextric, usando varios nodos
    de detección y procesamiento conectados en una red local.

    \item Analizar las latencias de las comunicaciones de los diferentes nodos en la red
    local (bien wifi, bien Ethernet), así como los tiempos de computación de cada etapa del
    procesamiento, factores decisivos para lograr un control efectivo en tiempo real.

    \item Diseñar y evaluar la detección y el seguimiento en tiempo real de la posición de
    cada coche en el circuito, empleando las cámaras necesarias, con o sin solapamiento
    entre sus campos de visión.

    \item Diseñar y evaluar el algoritmo de control adaptativo de la velocidad instantánea
    de cada coche en la arquitectura de control distribuida propuesta.

    \item Evaluar la arquitectura propuesta en condiciones reales de operación, midiendo el
    rendimiento, la latencia y la robustez del sistema.
\end{itemize}

\section{Propuesta de solución}
\label{sec:proposta}

Para implementar la arquitectura distribuida se emplea \gls{ros2}, un conjunto de
middleware, bibliotecas y herramientas de código abierto diseñado para el desarrollo de
aplicaciones robóticas complejas y escalables. Su adopción facilita el modelado del diseño
propuesto y garantiza la extensibilidad del sistema en el futuro; asimismo, simplifica la
gestión de las comunicaciones gracias a que incorpora protocolos de red optimizados y listos
para operar en tiempo real.

Sobre este middleware, el sistema se descompone en tres tipos de nodos que cooperan a través
de la red: un nodo cámara por cada cámara física, ejecutable en cualquier equipo de la red
local, que detecta la posición de los coches mediante segmentación por color; un nodo
controlador por cada rail que tengamos en circuito y por el que pueda circular un coche. 
Reconstruye la trayectoria del circuito durante una vuelta de calibración y ajusta en tiempo real 
la velocidad que hay que aplicar al carril en función de la posición junto con el comportamiento del vehículo; 
y un nodo puente que traduce las consignas de velocidad en
señales \gls{pwm} enviadas al Arduino que alimenta los carriles. El algoritmo de control
heredado de los trabajos anteriores se actualiza para soportar este procesamiento
distribuido, buscando maximizar el rendimiento sin comprometer la estabilidad.

Los tiempos de computación, la resolución de las imágenes, la latencia de las comunicaciones
y los retardos en el control \gls{pwm} son decisivos para lograr un control efectivo en
tiempo real, por lo que el sistema incorpora telemetría en todas sus etapas y herramientas
de grabación y análisis de las sesiones. La solución se completa con mecanismos de robustez
ante fallos y perturbaciones. También se realiza la publicación en tiempo real de información de interés: las imágenes procesadas por las cámaras, la posición de cada
coche y los tiempos por vuelta. Todo el sistema se despliega de forma reproducible mediante
contenedores Docker.

\section{Estructura de la memoria}
\label{sec:estrutura}

La memoria de este Trabajo de Fin de Grado se organiza en los siguientes capítulos:

\begin{itemize}
    \item \textbf{Capítulo 1. Introducción:} Se presenta el contexto general del proyecto,
    las motivaciones que lo impulsan, los objetivos planteados, la solución propuesta y la
    estructura que seguirá la memoria.

    \item \textbf{Capítulo 2. Tecnologías:} Se describen las tecnologías utilizadas durante
    el desarrollo del proyecto, tanto a nivel de hardware como de software.

    \item \textbf{Capítulo 3. Fundamentos teóricos y trabajos relacionados:} Se explican los
    conceptos técnicos relevantes y se revisan los trabajos previos que sirven como base o
    referencia para este proyecto.

    \item \textbf{Capítulo 4. Desarrollo del proyecto:} Se detallan los requisitos delo
    sistema, la metodología seguida, la planificación de tareas y los recursos empleados
    durante el proceso de desarrollo.

    \item \textbf{Capítulo 5. Diseño e implementación:} Se expone la arquitectura general
    del sistema distribuido y se describen los distintos nodos y componentes implementados,
    así como las decisiones tomadas en el diseño.

    \item \textbf{Capítulo 6. Pruebas:} Se presentan las pruebas realizadas para verificar
    el funcionamiento del sistema, junto con los resultados obtenidos y su análisis.

    \item \textbf{Capítulo 7. Conclusiones y trabajo futuro:} Se recogen las conclusiones
    del trabajo, se evalúa el grado de cumplimiento de los objetivos, se comentan las
    dificultades encontradas y se proponen posibles líneas de mejora o continuidad.
\end{itemize}
